\section{Conclusion}
\label{sec:conclusion}

We examined the dependency structure of Brazilian equities traded on B3 from 1998 to 2025 using an econophysics workflow. The main findings can be summarized along four dimensions.

First, the stylized-fact and correlation analysis shows that Brazilian equities display signatures commonly associated with complex financial systems: heavy-tailed return behaviour, volatility clustering, persistent absolute-return dependence and a moderate but positive average cross-correlation ($\bar{\rho}\approx0.24$). Within-sector correlations are larger than between-sector correlations on average, confirming that sectoral organization is reflected in the dependency matrix.

Second, Random Matrix Theory decomposes the correlation matrix. Five eigenvalues exceed the Mar\v{c}enko--Pastur upper bound, with the largest eigenvalue accounting for approximately 37.3\% of the trace of the correlation matrix. The leading component captures a market-mode component, while the remaining eigenvectors outside the random band suggest commodity, financial and utility-related group structures. The reconstruction checks indicate numerical accuracy within floating-point precision, and the group-mode matrix retains localized dependency patterns after the dominant market component is removed.

Third, financial network analysis through MST and PMFG representations shows that RMT filtering changes the structure of asset dependencies. The group-mode PMFG has lower mean edge correlation but higher clustering than the original PMFG, indicating stronger local clustering after market-mode removal. Hub rankings shift from networks dominated by broad market co-movement to networks where different assets occupy central positions after excluding the leading component. The PMFG retains 168 edges with triangles and 4-cliques, leaving local graph structures that are absent from the MST and useful for subsector-level analysis.

Fourth, the volatility forecasting experiment gives preliminary evidence that network topology can complement standard volatility predictors. Random Forest with network features obtains the best average QLIKE among the validation-selected models at the 5-day horizon, while Ridge regression with network features is strongest by QLIKE at the 20-day horizon. The CNN-1D model is competitive in MAE, RMSE and out-of-sample $R^2$, but the neural extension does not uniformly dominate across metrics. Feature-importance analysis shows that lagged volatility features dominate, while selected PMFG centralities contribute smaller but nonzero predictive information.

We view the net contribution of network features as incremental: they do not replace lagged-volatility dynamics, but they add structural information from the cross-sectional geometry of the market. Future work should extend this framework by constructing rolling RMT and network features for real-time forecasting, expanding the target asset set and evaluating richer temporal neural network and graph neural network architectures.
